\section{Franck--Condon 原理}
\label{ch:franck-condon}

前两章先从精确 Born--Huang 多通道理论降到单一绝热势能面，再把这一单势能面核哈密顿量特殊化为双原子分子的振动--转动问题。本章转而讨论两个不同电子势能面之间的光学跃迁。关键问题是：电子跃迁虽然很快，但核波函数并不会凭空消失；它在两个势能面上的投影决定了振动态谱线强度。单模位移谐振子只是最简单的可解模型，真实多原子分子还会出现频率改变和正规-模转动。
\subsection{从电子跃迁矩阵元到 Condon 近似}

在 BO 表示中，把初末态写成

\[
|g,\nu\rangle
=|\phi_g(\mathbf r;\mathbf Q)\rangle
|\chi_{g\nu}(\mathbf Q)\rangle,
\]

\[
|e,\nu'\rangle
=|\phi_e(\mathbf r;\mathbf Q)\rangle
|\chi_{e\nu'}(\mathbf Q)\rangle,
\]

其中 \(\mathbf Q\) 是质量加权正规坐标。电偶极矩阵元为

\[
M_{\nu'\nu}
=\int\dd\mathbf Q\,
\chi_{e\nu'}^*(\mathbf Q)
\mu_{eg}(\mathbf Q)
\chi_{g\nu}(\mathbf Q),
\]

电子跃迁偶极

\[
\mu_{eg}(\mathbf Q)
\equiv
\langle\phi_e(\mathbf Q)|\hat\mu|\phi_g(\mathbf Q)\rangle_{\mathbf r}.
\]

在参考构型 \(\mathbf Q=0\) 附近作 Taylor 展开：

\begin{equation}
\mu_{eg}(\mathbf Q)
=\mu_{eg}(0)
+\sum_\alpha
\left(\frac{\partial\mu_{eg}}{\partial Q_\alpha}\right)_0Q_\alpha
+O(Q^2).
\label{eq:herzberg-teller-expansion}
\end{equation}

Condon 近似是保留第一项：

\[
\mu_{eg}(\mathbf Q)\simeq\mu_{eg}(0).
\]

于是

\begin{equation}
M_{\nu'\nu}
\simeq
\mu_{eg}(0)
\langle\chi_{e\nu'}|\chi_{g\nu}\rangle,
\qquad
I_{\nu'\nu}\propto
|\langle\chi_{e\nu'}|\chi_{g\nu}\rangle|^2.
\label{eq:franck-condon-factor}
\end{equation}

若 \(\mu_{eg}(0)=0\) 是由对称性强制的，线性 Herzberg--Teller 项反而成为最低非零阶，此时“谱线强度只等于振动重叠平方”不再成立。因此 Franck--Condon 规则本身也有明确的截断条件。

Condon 近似把跃迁矩元简化为振动重叠问题，但重叠积分的计算仍需知道两个势能面的几何关系。最简单的情形是单模、等频率位移——此时 Franck--Condon 因子完全由 Huang--Rhys 因子 $S_{HR}$ 决定，且激发态振动真空在基态基底中恰为相干态。

\subsection{Huang--Rhys 因子}
先看一维、两条势能面曲率相同的模型：

\[
V_g(Q)=\frac12M\omega^2Q^2,
\]

\[
V_e(Q)=E_{eg}+\frac12M\omega^2(Q-\Delta Q)^2.
\]

定义零点涨落

\[
Q_{\rm zpf}=\sqrt{\frac{\hbar}{2M\omega}},
\]

以及无量纲位移

\begin{equation}
\lambda_{\rm FC}
\equiv\frac{\Delta Q}{2Q_{\rm zpf}}
=\sqrt{\frac{M\omega}{2\hbar}}\Delta Q,
\qquad
S_{\rm HR}\equiv\lambda_{\rm FC}^2.
\label{eq:huang-rhys}
\end{equation}

\(S_{\rm HR}\) 即 Huang--Rhys 因子。它可以理解为两个势阱平衡位置间位移相对于零点涨落宽度的平方，也等于经典重组能

\[
E_{\rm reorg}=\frac12M\omega^2(\Delta Q)^2
\]

与一个振动量子 \(\hbar\omega\) 的比值：

\[
S_{\rm HR}=\frac{E_{\rm reorg}}{\hbar\omega}.
\]

Huang--Rhys 因子给出了位移的标度，但 FC 因子的完整结构需用位移算符语言才能看清。以下将激发态振动真空写成基态位移算符作用于基态真空，揭示其相干态本质。

谐振子坐标写成

\[
Q=Q_{\rm zpf}(a+a^\dagger),
\qquad
P=-\ii\sqrt{\frac{M\hbar\omega}{2}}(a-a^\dagger).
\]

平移 \(Q\to Q-\Delta Q\) 的算符为

\[
\exp\left(-\frac{\ii}{\hbar}\Delta Q\,P\right)
=\exp[\lambda_{\rm FC}(a^\dagger-a)].
\]

因此定义位移算符

\begin{equation}
D(\lambda_{\rm FC})
=\exp[\lambda_{\rm FC}(a^\dagger-a)]
\label{eq:displacement-operator-fc}
\end{equation}

（这里 \(\lambda_{\rm FC}\) 取实数）。Baker--Campbell--Hausdorff 公式因为

\[
[\lambda_{\rm FC} a^\dagger,-\lambda_{\rm FC} a]=\lambda_{\rm FC}^2
\]

是 c 数，可精确写成

\begin{equation}
D(\lambda_{\rm FC})
=\ee^{-\lambda_{\rm FC}^2/2}
 \ee^{\lambda_{\rm FC} a^\dagger}
 \ee^{-\lambda_{\rm FC} a}.
\label{eq:displacement-normal-order}
\end{equation}

同时

\[
D^\dagger(\lambda_{\rm FC})aD(\lambda_{\rm FC})=a+\lambda_{\rm FC}.
\]

因此激发态振动真空相对于基态振动基底就是 displaced vacuum，也就是量子光学意义下的相干态。

\begin{proof}[解]
位移算符 \(D(\lambda)=\exp[\lambda(a^\dagger-a)]\)。坐标平移 \(Q\to Q-\Delta Q\) 由生成元 \(P\)（与 \(Q\) 共轭，\([Q,P]=\ii\hbar\)）的指数实现：\(T(\Delta Q)=\exp(-\ii\Delta Q\,P/\hbar)\)。由 BCH 公式 \(\ee^A B\ee^{-A}=B+[A,B]+\cdots\)，取 \(A=\ii\Delta Q\,P/\hbar\)：
\begin{equation*}
T^\dagger Q T=\ee^{\ii\Delta Q\,P/\hbar}\,Q\,\ee^{-\ii\Delta Q\,P/\hbar}=Q+[\tfrac{\ii\Delta Q\,P}{\hbar},Q]+\cdots
\end{equation*}
计算交换子：\([A,Q]=\frac{\ii\Delta Q}{\hbar}[P,Q]=\frac{\ii\Delta Q}{\hbar}(-\ii\hbar)=\Delta Q\)。高阶交换子 \([A,[A,Q]]=[A,\Delta Q]=0\) 截断。因此 \(T^\dagger Q T=Q+\Delta Q\)。代入 \(P=-\ii\sqrt{M\hbar\omega/2}(a-a^\dagger)\)：
\begin{equation*}
T=\exp\!\left(-\frac{\ii\Delta Q}{\hbar}\cdot(-\ii)\sqrt{\frac{M\hbar\omega}{2}}(a-a^\dagger)\right)
=\exp\!\left(-\frac{\Delta Q}{\hbar}\sqrt{\frac{M\hbar\omega}{2}}(a-a^\dagger)\right)
=\exp[\lambda_{\rm FC}(a^\dagger-a)],
\end{equation*}
其中 \(\lambda_{\rm FC}=\Delta Q\sqrt{M\omega/(2\hbar)}=\Delta Q/(2Q_{\rm zpf})\)，即\eqnrefs{eq:huang-rhys}的定义。

令 \(A=\lambda a^\dagger\)，\(B=-\lambda a\)，\([A,B]=\lambda^2\) 为 c 数，BCH 公式 \(\ee^{A+B}=\ee^A\ee^B\ee^{-[A,B]/2}\) 给出
\begin{equation*}
D(\lambda)=\ee^{\lambda(a^\dagger-a)}=\ee^{\lambda a^\dagger}\ee^{-\lambda a}\ee^{-\lambda^2/2},
\end{equation*}
即\eqnrefs{eq:displacement-normal-order}。

由 \([a,a^\dagger]=1\)，\([a,\lambda a^\dagger]=\lambda\)，\([a,-\lambda a]=0\)。BCH 公式：
\begin{equation*}
D^\dagger(\lambda)aD(\lambda)=\ee^{-\lambda(a^\dagger-a)}\,a\,\ee^{\lambda(a^\dagger-a)}
=a+[\lambda(a^\dagger-a),a]+\frac{1}{2}[\lambda(a^\dagger-a),[\lambda(a^\dagger-a),a]]+\cdots
\end{equation*}
一阶：\([\lambda(a^\dagger-a),a]=\lambda[a^\dagger,a]-\lambda[a,a]=-\lambda\)。二阶：\([\lambda(a^\dagger-a),-\lambda]=0\) 截断。因此
\begin{equation*}
D^\dagger a D=a+\lambda.
\end{equation*}
即 \(a(D|0\rangle)=\lambda(D|0\rangle)\)。位移算符 \(D(\lambda_{\rm FC})\) 把基态振动真空平移到激发态振动真空，\(\lambda_{\rm FC}=\Delta Q/(2Q_{\rm zpf})\) 即两势阱平衡位置之差以零点宽度为单位，FC 因子由 Huang--Rhys 因子 \(S_{\rm HR}=\lambda_{\rm FC}^2\) 决定。综合以上，位移算符的关键代数性质为
\begin{equation*}
D^\dagger(\lambda)\,a\,D(\lambda)=a+\lambda,\qquad D^\dagger(\lambda)\,a^\dagger\,D(\lambda)=a^\dagger+\lambda^*.
\end{equation*}

\end{proof}

有了真空-真空重叠，一般 $m\to n$ 跃迁的 FC 因子只需在位移算符两侧插入升降算符。以下用生成函数方法一次性获得全部 FC 因子的封闭表达式。

设初态为基态势阱的振动真空。由于

\[
\ee^{-\lambda_{\rm FC} a}|0\rangle=|0\rangle,
\]

有

\[
D(\lambda_{\rm FC})|0\rangle
=\ee^{-\lambda_{\rm FC}^2/2}
\sum_{n=0}^{\infty}
\frac{\lambda_{\rm FC}^n}{n!}(a^\dagger)^n|0\rangle.
\]

利用

\[
(a^\dagger)^n|0\rangle=\sqrt{n!}|n\rangle,
\]

得到

\begin{equation}
\langle n|D(\lambda_{\rm FC})|0\rangle
=\ee^{-\lambda_{\rm FC}^2/2}\frac{\lambda_{\rm FC}^n}{\sqrt{n!}}.
\label{eq:fc-amplitude-ground}
\end{equation}

因此

\begin{equation}
|\langle n_e|0_g\rangle|^2
=\ee^{-S_{\rm HR}}
\frac{S_{\rm HR}^n}{n!}.
\label{eq:fc-poisson}
\end{equation}

不同 Huang--Rhys 因子下的分布见\figref{fig:franck-condon}。

\begin{figure}[tbp]
\centering
\includegraphics[width=.83\textwidth]{franck_condon.pdf}
\caption{等频率位移谐振子模型中的 Franck--Condon 因子。Huang--Rhys 因子 \(S_{\rm HR}\) 同时决定 Poisson 分布的均值和方差。}
\label{fig:franck-condon}
\end{figure}

\eqnrefs{eq:fc-poisson} 的均值和方差都是 \(S_{\rm HR}\)。因此 \(S_{\rm HR}\ll1\) 时零声子线占主导；\(S_{\rm HR}\gg1\) 时强度分布中心移向 \(n\sim S_{\rm HR}\)，出现明显振动 progression。

\begin{proof}[解]
对任意数态，使用 \eqnrefs{eq:displacement-normal-order}：
\begin{equation*}
\langle m|D(\lambda_{\rm FC})|n\rangle
=\ee^{-\lambda_{\rm FC}^2/2}
\langle m|\ee^{\lambda_{\rm FC} a^\dagger}\ee^{-\lambda_{\rm FC} a}|n\rangle.
\end{equation*}
先展开右侧的 \(\ee^{-\lambda_{\rm FC} a}\)，利用 \(a^j|n\rangle=\sqrt{n!/(n-j)!}\,|n-j\rangle\) 且 \(a^j|n\rangle=0\) 对 \(j>n\)，级数自动截断：
\begin{equation*}
\ee^{-\lambda_{\rm FC} a}|n\rangle
=\sum_{j=0}^{\infty}\frac{(-\lambda_{\rm FC})^j}{j!}a^j|n\rangle
=\sum_{j=0}^{n}\frac{(-\lambda_{\rm FC})^j}{j!}
\sqrt{\frac{n!}{(n-j)!}}\,|n-j\rangle.
\end{equation*}
再展开左侧的 \(\ee^{\lambda_{\rm FC} a^\dagger}\) 并利用数态正交性 \(\langle m|a^{\dagger i}|n-j\rangle=\sqrt{m!/(m-i)!}\,\delta_{m-i,n-j}\)，得到
\begin{align*}
\langle m|D(\lambda_{\rm FC})|n\rangle
&=\ee^{-\lambda_{\rm FC}^2/2}
\sum_{j=0}^{n}\frac{(-\lambda_{\rm FC})^j}{j!}
\sqrt{\frac{n!}{(n-j)!}}
\sum_{i=0}^{\infty}\frac{\lambda_{\rm FC}^i}{i!}
\sqrt{\frac{m!}{(m-i)!}}\,\delta_{m-i,n-j}\\
&=\ee^{-\lambda_{\rm FC}^2/2}\sqrt{n!\,m!}
\sum_{j=0}^{n}\frac{(-1)^j\lambda_{\rm FC}^{m-n+2j}}{(m-n+j)!\,j!\,(n-j)!},
\end{align*}
其中 Kronecker 记号固定 \(i=m-n+j\)，故求和指标仍为 \(j\)，且要求 \(m\ge n\)。把 \(\lambda_{\rm FC}^{m-n+2j}=\lambda_{\rm FC}^{m-n}\lambda_{\rm FC}^{2j}\) 拆开并提出求和号外，引入关联 Laguerre 多项式
\begin{equation*}
L_n^{(\alpha)}(x)
=\sum_{j=0}^{n}\frac{(-x)^j}{j!}\,
\frac{\Gamma(n+\alpha+1)}{\Gamma(\alpha+j+1)\,(n-j)!},
\end{equation*}
取 \(\alpha=m-n\) 时 \(\Gamma(m+1)=m!\)、\(\Gamma(m-n+j+1)=(m-n+j)!\)，于是求和正好满足
\begin{equation*}
\sum_{j=0}^{n}\frac{(-1)^j\lambda_{\rm FC}^{m-n+2j}}{(m-n+j)!\,j!\,(n-j)!}
=\frac{\lambda_{\rm FC}^{m-n}}{m!}
L_n^{(m-n)}(\lambda_{\rm FC}^2),
\end{equation*}
代回前式并约去 \(m!\)，
\begin{align*}
\langle m|D(\lambda_{\rm FC})|n\rangle
&=\ee^{-\lambda_{\rm FC}^2/2}\sqrt{n!\,m!}\,
\frac{\lambda_{\rm FC}^{m-n}}{m!}
L_n^{(m-n)}(\lambda_{\rm FC}^2)\\
&=\ee^{-\lambda_{\rm FC}^2/2}
\sqrt{\frac{n!}{m!}}\,\lambda_{\rm FC}^{m-n}
L_n^{(m-n)}(\lambda_{\rm FC}^2).
\end{align*}
因此对于 \(m\ge n\)，
\begin{equation*}
\langle m|D(\lambda_{\rm FC})|n\rangle
=\ee^{-\lambda_{\rm FC}^2/2}
\sqrt{\frac{n!}{m!}}
\lambda_{\rm FC}^{m-n}
L_n^{(m-n)}(\lambda_{\rm FC}^2),
\end{equation*}
而 \(n\ge m\) 时由厄米共轭 \(\langle m|D|n\rangle=\langle n|D^\dagger|m\rangle^*\) 或直接交换 \(m,n\) 得到相应表达。这里 \(L_n^{(\alpha)}\) 是广义 Laguerre 多项式。\(n=0\) 时 \(L_0^{(m)}=1\)，立刻退化为 \eqnrefs{eq:fc-amplitude-ground}。
\end{proof}

等频率位移只涉及位移算符 $D(\alpha)$。但实际分子的两个电子态势能面通常频率不同、正规模方向也不一致，分别引入压缩和 Duschinsky 旋转。本节处理这一一般情形。

\subsection{Duschinsky--Doktorov 变换}
如果两电子态的振动频率不同，

\[
\omega_g\neq\omega_e,
\]

即使平衡位置相同，两组升降算符也不同。用同一组物理 \(Q,P\) 写

\[
a_g
=\sqrt{\frac{M\omega_g}{2\hbar}}Q
+\frac{\ii}{\sqrt{2M\hbar\omega_g}}P,
\]

\[
a_e
=\sqrt{\frac{M\omega_e}{2\hbar}}Q
+\frac{\ii}{\sqrt{2M\hbar\omega_e}}P.
\]

消去 \(Q,P\)，得到 Bogoliubov 变换

\[
a_e
=u a_g+v a_g^\dagger,
\]

\[
u=\frac12\left(
\sqrt{\frac{\omega_e}{\omega_g}}
+\sqrt{\frac{\omega_g}{\omega_e}}
\right),
\qquad
v=\frac12\left(
\sqrt{\frac{\omega_e}{\omega_g}}
-\sqrt{\frac{\omega_g}{\omega_e}}
\right),
\]

满足 \(u^2-v^2=1\)。令

\[
r=\frac12\ln\frac{\omega_e}{\omega_g},
\]

则 \(u=\cosh r,v=\sinh r\)。所以“频率改变对应压缩”不是类比，而是两套谐振子正则坐标之间的精确 Bogoliubov 关系。

频率改变只引入对角压缩。多原子分子中两个电子态的正规坐标方向一般不同，还需一个正交旋转——Duschinsky 变换——来描述模式混合。

多原子分子的两个电子态一般具有不同 Hessian，因此正规模方向会旋转。采用质量加权正规坐标，最常见的线性关系写成

\begin{equation}
\mathbf Q_e
=\mathbf J\mathbf Q_g+\mathbf K,
\label{eq:duschinsky}
\end{equation}

其中 \(\mathbf J\) 是 Duschinsky 转动，\(\mathbf K\) 是平衡构型位移。若两组质量加权坐标都正交归一，则 \(\mathbf J\) 为正交矩阵，并且为保持正则结构

\[
\mathbf P_e=\mathbf J\mathbf P_g.
\]

令 \(\boldsymbol\Omega_g\)、\(\boldsymbol\Omega_e\) 分别是由各振动频率组成的对角矩阵，定义多模升降算符

\[
\mathbf a_s
=\frac1{\sqrt{2\hbar}}
\left(
\boldsymbol\Omega_s^{1/2}\mathbf Q_s
+\ii\boldsymbol\Omega_s^{-1/2}\mathbf P_s
\right),
\qquad s=g,e.
\]

把 \eqnrefs{eq:duschinsky} 代入后得到一般高斯仿射变换

\begin{equation}
\mathbf a_e
=\mathbf X\mathbf a_g
+\mathbf Y\mathbf a_g^\dagger
+\mathbf z,
\label{eq:doktorov-bogoliubov}
\end{equation}

其中

\[
\mathbf X
=\frac12
\left(
\boldsymbol\Omega_e^{1/2}\mathbf J\boldsymbol\Omega_g^{-1/2}
+
\boldsymbol\Omega_e^{-1/2}\mathbf J\boldsymbol\Omega_g^{1/2}
\right),
\]

\[
\mathbf Y
=\frac12
\left(
\boldsymbol\Omega_e^{1/2}\mathbf J\boldsymbol\Omega_g^{-1/2}
-
\boldsymbol\Omega_e^{-1/2}\mathbf J\boldsymbol\Omega_g^{1/2}
\right),
\]

\[
\mathbf z
=\frac1{\sqrt{2\hbar}}
\boldsymbol\Omega_e^{1/2}\mathbf K.
\]

因此多模 Franck--Condon 变换不是简单的 \(D S U\) 口号，而是一个保持 bosonic 对易子的多模 Bogoliubov 变换加位移。\(\mathbf J\) 负责模式混合，\(\mathbf X,\mathbf Y\) 中的频率比负责压缩，\(\mathbf z\) 负责位移。它们合起来就是常说的 Doktorov 变换。

实际大分子的 Franck--Condon 因子因而是多维高斯波函数之间的重叠问题；模式旋转和压缩会让不同简正模之间的量子数不再逐模独立守恒。

以上均在零温（初态为振动基态）下讨论。有限温度下初态为热态，FC 因子需对热平均求和，用时间关联函数可以给出紧凑的表达。

\subsection{有限温度谱与时间关联函数}

在有限温度下，初始振动态不再是单个 $|0_g\rangle$，而是热态
\[
\rho_g^{\rm th}=\frac{\ee^{-\beta H_g}}{Z_g},
\qquad \beta=(k_BT)^{-1}.
\]
Condon 近似下，吸收谱可以避免逐个枚举所有初末振动态，而写成偶极关联函数的 Fourier 变换
\begin{equation}
I(\omega)\propto
\int_{-\infty}^{\infty}\dd t\,
\ee^{\ii\omega t}
C_{eg}(t),
\qquad
C_{eg}(t)=\operatorname{Tr}\!\left[
ho_g^{\rm th}
\ee^{\ii H_gt/\hbar}\ee^{-\ii H_et/\hbar}\right].
\label{eq:franck-condon-thermal-correlation}
\end{equation}
对一般二次 $H_g,H_e$，\eqnrefs{eq:doktorov-bogoliubov}保证这个迹仍是高斯迹，因此可以由协方差矩阵和位移向量闭合计算。零温、等频率、单模位移模型只是在该一般生成函数中取 $T\to0$、$\mathbf J=I$、$\mathbf Y=0$ 后得到的 Poisson 特例。

\begin{proof}[解]
离散振动态的吸收强度由 Fermi 黄金法则给出
\begin{equation*}
I(\omega)\propto
\sum_{m,n} p_n\,
|\langle e,m|\hat\mu|g,n\rangle|^2\,
\delta\!\left(\omega-\frac{E_m^e-E_n^g}{\hbar}\right),
\end{equation*}
其中 $p_n=\ee^{-\beta E_n^g}/Z_g$ 为初态 Boltzmann 权重，$|g,n\rangle$、$|e,m\rangle$ 分别为电子基态与激发态势能面上的振动态。记 $\omega_{mn}=(E_m^e-E_n^g)/\hbar$，将矩阵元模平方展开，
\begin{equation*}
|\langle e,m|\hat\mu|g,n\rangle|^2
= \langle e,m|\hat\mu|g,n\rangle\,
 \langle g,n|\hat\mu^\dagger|e,m\rangle.
\end{equation*}
利用 $\delta$ 函数的 Fourier 表示
\begin{equation*}
\delta(\omega-\omega_{mn})
= \frac1{2\pi}\int_{-\infty}^{\infty}\dd t\,
 \ee^{\ii(\omega-\omega_{mn})t},
\end{equation*}
代回强度公式并交换求和与积分，
\begin{equation*}
I(\omega)\propto
\frac1{2\pi}\int_{-\infty}^{\infty}\dd t\,
\ee^{\ii\omega t}
\sum_{m,n} p_n\,
\langle e,m|\hat\mu|g,n\rangle
\langle g,n|\hat\mu^\dagger|e,m\rangle
\ee^{-\ii\omega_{mn}t}.
\end{equation*}
由 $H_g|g,n\rangle=E_n^g|g,n\rangle$、$H_e|e,m\rangle=E_m^e|e,m\rangle$，于是
\begin{equation*}
\ee^{-\ii\omega_{mn}t}
= \langle e,m|\ee^{-\ii H_e t/\hbar}|e,m\rangle\,
 \langle g,n|\ee^{+\ii H_g t/\hbar}|g,n\rangle.
\end{equation*}
在 Condon 近似下电子偶极与核坐标无关，$\hat\mu=\boldsymbol\mu_{eg}\otimes\hat{\mathbf 1}_{\rm vib}$，可提出电子因子。对 $m$ 与 $n$ 分别插入完备关系 $\sum_m|e,m\rangle\langle e,m|=\mathbf 1$、$\sum_n|g,n\rangle\langle g,n|=\mathbf 1$，双重求和坍缩为
\begin{equation*}
\sum_n p_n\,
\langle g,n|\ee^{+\ii H_g t/\hbar}\ee^{-\ii H_e t/\hbar}|g,n\rangle
= \operatorname{Tr}\!\left[
 \rho_g^{\rm th}\,
 \ee^{+\ii H_g t/\hbar}\ee^{-\ii H_e t/\hbar}\right],
\end{equation*}
其中 $\rho_g^{\rm th}=\sum_n p_n|g,n\rangle\langle g,n|=\ee^{-\beta H_g}/Z_g$。因此
\begin{equation*}
I(\omega)\propto
\int_{-\infty}^{\infty}\dd t\,
\ee^{\ii\omega t}\,
C_{eg}(t),
\qquad
C_{eg}(t)=\operatorname{Tr}\!\left[
\rho_g^{\rm th}\,
\ee^{+\ii H_g t/\hbar}\ee^{-\ii H_e t/\hbar}\right],
\end{equation*}
即\eqnrefs{eq:franck-condon-thermal-correlation}。大分子振动谱更自然的母对象是高斯关联函数，而非数量指数增长的逐态 Franck--Condon 表。
\end{proof}

\subsection{Franck--Condon 因子的实验数值与典型体系}\label{sec:fc-experiments}

Huang--Rhys 因子、Duschinsky 旋转和有限温度关联函数分别控制位移、模式混合与热布居。代入具体分子的振动频率和平衡构型差以后，它们直接给出谱带包络和各振动支线的相对强度。

单模位移谐振子的 Huang--Rhys 因子 \(S_{\rm HR}=(\Delta Q)^2/(2\ell_v^2)\) 决定振动电子谱线强度分布：\(I_n=e^{-S_{\rm HR}}S_{\rm HR}^{\,n}/n!\)（Poisson 分布）。\(S_{\rm HR}\) 的物理意义是基态到末态振动波包中心的位移（以零点宽度为单位）。

\begin{center}
\begin{tabular}{lcccc}
\hline
体系 & 跃迁 & \(S\) & Stokes 位移 (cm\(^{-1}\)) & \(\lambda_{\rm FC}^2\) \\
\hline
H\(_2\) (\(B\,{}^1\Sigma\leftarrow X\,{}^1\Sigma\)) & Lyman & 0.09 & 1700 & 0.18 \\
I\(_2\) (\(B\,{}^3\Pi\leftarrow X\,{}^1\Sigma\)) & 可见 & 1.5 & 240 & 3.0 \\
F-center (NaCl) & \(1s\to2p\) & 3.0 & 1800 & 6.0 \\
聚芴 (有机半导体) & \(S_1\to S_0\) & 2.0 & 2400 & 4.0 \\
CdSe 量子点 & 带边 & 0.5 & 200 & 1.0 \\
\hline
\end{tabular}
\end{center}

\begin{exbox}[\(I_2\) 的振动电子进展]
I\(_2\) 的 \(B\leftarrow X\) 跃迁 \(S\approx1.5\)，振动电子峰强分布 \(I_n=e^{-1.5}1.5^n/n!\)：\(I_0=0.22\)，\(I_1=0.34\)，\(I_2=0.25\)，\(I_3=0.13\)，\(I_4=0.05\)。最强峰在 \(n\approx S=1.5\) 附近，进展宽度 \(\Delta n\approx\sqrt{S}\approx1.2\)。对比 H\(_2\) 的 \(S=0.09\)：\(I_0=0.91\)，\(I_1=0.08\)——几乎只有 0-0 带，因为 H\(_2\) 两个电子态的平衡键长几乎相同。
\end{exbox}

Huang--Rhys 因子决定振动电子包络的形状，但光谱中哪些跃迁实际出现还受选择定则约束。大 $S$ 时反射原理将 FC 因子与连续谱阈联系起来。

\begin{exbox}[反射原理：\(I_2\) 吸收谱的连续阈]
当 \(S\gg1\) 且末态势能面显著位移时，FC 因子在 \(v'\approx S\) 处达到峰值，对应经典转折点。\(I_2\) 的 \(B\) 态势能在 \(R\approx3.0\,\text{\AA}\) 处有势阱，而基态 \(R_e=2.67\,\text{\AA}\)。当激发光能量超过 \(B\) 态离解限时，FC 因子给出连续吸收谱——反射原理：核波函数在转折点反射，吸收谱形 \(\sigma(\nu)\propto|\chi_g(R_t)|^2/|V_e'(R_t)-V_g'(R_t)|\)，其中 \(R_t\) 满足 \(V_e(R_t)-V_g(R_t)=h\nu\)。
\end{exbox}

以上实验体系均为单模或频率一致情形。一般多原子分子中 Duschinsky 旋转使模式混合，FC 因子成为多维重叠积分。

当两个电子态的正规-模方向不同时，Duschinsky 变换 \(Q'=UQ+\Delta Q\) 引入模式混合。单模 FC 因子推广为多维重叠积分。

\begin{exbox}[\(CO_2\) 的 Duschinsky 效应]
CO\(_2\) 在 \(S_1\leftarrow S_0\) 跃迁中，弯曲模 \(\nu_2\) 在激发态变为非对称伸缩，Duschinsky 矩阵 \(U\) 非对角。忽略 Duschinsky 转动会使振动电子峰强误差达 30\%。正确处理需计算多维 FC 积分，常用 Doktorov 算法（递推生成）或热关联函数方法。
\end{exbox}

\begin{proof}[解]

设初态 \(g\) 有 \(N\) 个简正模 \(\{Q_i\}\)，频率 \(\{\omega_i\}\)；末态 \(e\) 有 \(N\) 个模 \(\{Q'_j\}\)，频率 \(\{\omega'_j\}\)。Duschinsky 变换
\begin{equation}
\mathbf Q' = U\,\mathbf Q + \mathbf d,
\label{eq:duschinsky-general}
\end{equation}
其中 \(U\in\mathrm{O}(N)\) 为正交矩阵（模式混合），\(\mathbf d\) 为位移向量。多模 FC 因子是 \(N\) 维重叠积分
\begin{equation}
F_{\mathbf n,\mathbf m}
=\int\!\dd^NQ\;\chi_{g,\mathbf n}(\mathbf Q)\,\chi_{e,\mathbf m}(U\mathbf Q+\mathbf d),
\label{eq:multimode-fc-integral}
\end{equation}


引入热关联函数
\begin{equation}
C(t)=\frac{1}{Z_g}\Tr\!\left[\ee^{-\beta H_g}\,\ee^{\ii H_e t/\hbar}\,\ee^{-\ii H_g t/\hbar}\right],
\qquad
Z_g=\Tr\,\ee^{-\beta H_g},
\label{eq:thermal-correlation}
\end{equation}
其中 \(H_g=\sum_i\hbar\omega_i(a_i^\dagger a_i+1/2)\)，\(H_e=\sum_j\hbar\omega'_j(b_j^\dagger b_j+1/2)\)，且末态算符通过 Duschinsky 变换与初态关联。吸收谱强度
\begin{equation*}
I(\omega)\propto\int_{-\infty}^{+\infty}\!\dd t\;\ee^{-\ii\omega t}\,C(t),
\end{equation*}


用相干态 \(|\boldsymbol\alpha\rangle\) 表示，\(\boldsymbol\alpha\in\CC^N\)。路径积分中 \(C(t)\) 化为 \(2N\) 维高斯积分。定义复对称矩阵
\begin{equation}
\Sigma(t)_{ij}
=\frac{\delta_{ij}\,\ee^{-\ii\omega_i t}}{1-\ee^{-\beta\hbar\omega_i-\ii\omega_i t}},
\label{eq:coherent-state-propagator}
\end{equation}
以及 Duschinsky 诱导的频率比矩阵 \(W_{ij}=\sqrt{\omega'_j/\omega_i}\,U_{ij}\)。完成高斯积分后
\begin{equation}
C(t)=\frac{\det{}^{1/2}\!\bigl[\mathbb I-W^\dagger\Sigma(t)W\Sigma(-t)\bigr]}
{\prod_i\bigl(1-\ee^{-\beta\hbar\omega_i-\ii\omega_i t}\bigr)^{1/2}}
\,\exp\!\Bigl[-\frac12\mathbf d^\dagger\Lambda(t)\mathbf d\Bigr],
\label{eq:multimode-fc-generating-function}
\end{equation}
其中 \(\Lambda(t)\) 为由 \(W\)、\(\Sigma(t)\) 构成的复对称矩阵。行列式项来自模式混合，指数项来自位移。

正则对易关系 \([Q_i,P_j]=\ii\hbar\delta_{ij}\) 在 Duschinsky 变换下保持，要求
\begin{equation}
W^\dagger\Omega_e W=\Omega_g,
\qquad
\Omega_{g,ij}=\ii\hbar\omega_i\delta_{ij},
\qquad
\Omega_{e,ij}=\ii\hbar\omega'_j\delta_{ij},
\label{eq:duschinsky-symplectic}
\end{equation}


取单模极限 \(N=1\)、\(U=1\)、\(\mathbf d=d\)。此时 \(W=\sqrt{\omega'/\omega}\)，若进一步 \(\omega'=\omega\)（无频率改变），\eqnrefs{eq:multimode-fc-generating-function} 简化为
\begin{equation*}
C(t)=\exp\!\Bigl[-S\bigl(1-\ee^{-\ii\omega t}\bigr)\Bigr],
\qquad
S=\frac{\omega d^2}{2\hbar},
\end{equation*}
展开 \(\ee^{S\ee^{-\ii\omega t}}=\sum_k(S^k/k!)\ee^{-\ii k\omega t}\)，Fourier 变换给出
\begin{equation*}
I_k\propto\ee^{-S}\frac{S^k}{k!},
\end{equation*}
即\eqnrefs{eq:fc-poisson} 的 Poisson 分布。

\end{proof}

多模 Franck--Condon 因子还可直接由 Duschinsky 变换构造。

Duschinsky 变换。电子跃迁改变势能面，初态正则坐标 \(\boldsymbol Q'\) 与末态 \(\boldsymbol Q\) 通过 Duschinsky 变换关联：\(\boldsymbol Q'=D\boldsymbol Q+\boldsymbol d\)（\(D\) 正交矩阵描述模混合，\(\boldsymbol d\) 位移向量）。\(D\) 的非对角元表示初末态正则模不同（频率或力常数矩阵不同），\(\boldsymbol d\) 表示平衡位置偏移。单模 FC（\(D=I\)）只涉及位移，多模 FC 需要处理 \(D\neq I\) 的模混合。

多模重叠最好由一个生成函数统一保存。记
\[
\mathcal F_{\boldsymbol v',\boldsymbol v}
=\langle\boldsymbol v'|D,\boldsymbol d|\boldsymbol v\rangle ,
\]
并定义
\begin{equation}
\begin{aligned}
G(\boldsymbol z,\boldsymbol z')
=\sum_{\boldsymbol v,\boldsymbol v'}
&\mathcal F_{\boldsymbol v',\boldsymbol v}
\prod_{i=1}^{N}\frac{z_i^{v_i}}{\sqrt{v_i!}}
\prod_{j=1}^{N}\frac{{z'_j}^{v'_j}}{\sqrt{v'_j!}} .
\end{aligned}
\label{eq:multimode-fc-amplitude-generator}
\end{equation}
两个势能面都是二次型时，所有积分仍是高斯积分，因此生成函数必可写成
\begin{equation}
G=\exp\!\left[
\tfrac12\boldsymbol z^{\mathsf T}A\boldsymbol z
+\tfrac12{\boldsymbol z'}^{\mathsf T}B\boldsymbol z'
+\boldsymbol z^{\mathsf T}C\boldsymbol z'
+\boldsymbol l^{\mathsf T}\boldsymbol z
+{\boldsymbol l'}^{\mathsf T}\boldsymbol z'
\right],
\label{eq:doktorov-gaussian-generator}
\end{equation}
其中 \(A,B,C,\boldsymbol l,\boldsymbol l'\) 由 Duschinsky 矩阵、位移和两组频率决定。对 \(\boldsymbol z,\boldsymbol z'\) 求导并在原点取值，就逐项取回所需的 Franck--Condon 振幅。

单模位移常用 Huang--Rhys 因子
\[
S_{\rm HR}=\frac{M\omega d^2}{2\hbar}
\]
度量。若 \(v'\ge v\)，重叠可写成
\begin{equation}
\langle v'|v\rangle
=\sqrt{\frac{v!}{v'!}}\,
\ee^{-S_{\rm HR}/2}
S_{\rm HR}^{(v'-v)/2}
L_v^{(v'-v)}(S_{\rm HR}),
\label{eq:huang-rhys-laguerre-overlap}
\end{equation}
\(v'<v\) 时交换上下指标。\(S_{\rm HR}\ll1\) 时零声子线占主导；\(S_{\rm HR}\gg1\) 时强度移到 \(v'\) 约为 \(S_{\rm HR}\) 的振动边带。这个无量纲量因而把势能面的几何位移直接翻译成电子--声子耦合强度。

直接展开 \eqnrefs{eq:multimode-fc-amplitude-generator} 会随模式数迅速膨胀。Doktorov 递推把高阶振动态降到相邻低阶态，典型一步具有形式
\begin{equation}
\begin{aligned}
\sqrt{v_i}\,
\mathcal F_{\boldsymbol v',\,\boldsymbol v+\boldsymbol e_i}
={}&\sum_j M_{ij}\sqrt{v'_j+1}\,
\mathcal F_{\boldsymbol v'+\boldsymbol e'_j,\,\boldsymbol v}\\
&+\sum_j N_{ij}\sqrt{v'_j}\,
\mathcal F_{\boldsymbol v'-\boldsymbol e'_j,\,\boldsymbol v}
+P_i\mathcal F_{\boldsymbol v',\,\boldsymbol v}.
\end{aligned}
\label{eq:doktorov-fc-recurrence}
\end{equation}
系数 \(M,N,P\) 同样由 \(D,\boldsymbol d\) 和频率决定。递推避免枚举生成函数的全部 Taylor 项，但计算量仍随所保留的振动态数
\(\prod_i(v_i^{\max}+1)\) 增长；因此“多项式”只针对给定截断网格，并不意味着百模强激发问题自动变得廉价。
